% !TeX root = from-differential-to-algebraic-geometry.tex


\subsection{Almost complex structures}

Let \(X\) be a complex manifold of complex dimension \(n\).
The tangent bundle \(TX\) is a real vector bundle of rank \(2n\).
There is a natural endomorphism \(J: TX\to TX\) induced by the complex structure of \(X\), i.e., for every point \(p\in X\), \(J_p: T_pX\to T_pX\) is the multiplication by \(\im\).
We have \(J^2 = -\id\).

\begin{question}\label{qs:existence-and-uniqueness-of-complex-structure}
	Given a topological space \(M\) that is Hausdorff and second countable, when does it admit a complex structure?
	Is such a complex structure unique?
\end{question}

For complex dimension \(1\), the answer is positive and well-known.
For higher dimensions, the answer is negative in general.
In particular, does the \(6\)-sphere \(S^6\) admit a complex structure?
This is a famous open problem in complex geometry.

\begin{question}\label{qs:complex_structure_on_S6}
	Does the \(6\)-sphere \(S^6\) admit a complex structure?
\end{question}

\begin{definition}\label{def:almost_complex_structure}
	Let \(M\) be a smooth manifold of real dimension \(2n\).
	An \emph{almost complex structure} on \(M\) is a smooth endomorphism \(J: TM\to TM\) such that \(J^2 = -\id\).
	The pair \((M,J)\) is called an \emph{almost complex manifold}.
\end{definition}

% \begin{definition}\label{def:J_holomorphic_function}
%     Let \((M,J)\) be an almost complex manifold.
%     A smooth function \(f: M\to \bbC\) is called \emph{\(J\)-holomorphic} if
%     \[ \upd f \circ J = \im \cdot \upd f. \]
% \end{definition}

\begin{question}\label{qs:existence_of_almost_complex_structure}
	Given a smooth manifold \(M\) of real dimension \(2n\), when does it admit an almost complex structure?
	Is such an almost complex structure unique?
\end{question}

Giving an almost complex structure \(J\) on a smooth manifold \(M\) is equivalent to giving the tangent bundle \(TM\) the structure of a complex vector bundle.
Hence the existence of almost complex structures is a purely topological problem.
Note that to find a complex structure on \(M\) needs to solve some non-linear partial differential equations, which is much harder.

\begin{example}\label{eg:almost_complex_structure_on_S6}
	The \(6\)-sphere \(S^6\) admits an almost complex structure.
	In fact, \(S^6\) can be identified with the unit sphere in the imaginary octonions \(\Image \bbO\) (see \cref{rmk:fundamental_facts_about_octonion_for_almost_complex_structure_on_S6}).
	Denote by \(m(x,y)\) the octonionic multiplication of \(x,y\in \bbO\).
	For every point \(p\in S^6\), the tangent space \(T_pS^6\) can be identified with the orthogonal complement of \(\bbR p\) in \(\Image \bbO\).
	Define \(J_p: T_pS^6 \to T_pS^6\) by \(J_p(v) = m(p,v)\).
	Then \(J_p^2(v) = p(pv) = -v\) for every \(v\in T_pS^6\).
	Thus we get an almost complex structure on \(S^6\).
\end{example}
\begin{remark}\label{rmk:fundamental_facts_about_octonion_for_almost_complex_structure_on_S6}
	Recall some fundamental facts about the octonions \(\bbO\):
	\begin{enumerate}
		\item \(\bbO\) is an \(8\)-dimensional normed vector space over \(\bbR\) with an orthogonal basis \(\{1\}\cup \{e_i|i=1,\ldots,7\}\).
		      The subspace spanned by \(\{e_i\}\) is called the space of imaginary octonions and denoted by \(\Image \bbO\).
		\item The multiplication \(m: \bbO \times \bbO \to \bbO\) is a bilinear map and satisfies the distributive law and the norm multiplicative law \(\|xy\| = \|x\|\|y\|\) for all \(x,y\in \bbO\).
		      It is given by the following Fano plane \(\bbP^2(\bbF_2)\):
		      \begin{center}
			      \usetikzlibrary{calc, decorations.markings}
			      \begin{tikzpicture}[scale=1.5, >=stealth,
					      midarrow/.style={decoration={markings, mark=at position 0.5 with {\arrow{>}}}, postaction={decorate}}]
				      % Define the vertices of the triangle
				      \coordinate (A) at (0, 2);              %e6
				      \coordinate (B) at (-1.732, -1);        %e3
				      \coordinate (C) at (1.732, -1);         %e5
				      % Define the center and midpoints
				      \coordinate (O) at (0, 0);              %e7
				      \coordinate (D) at (0, -1);             %e2
				      \coordinate (E) at (0.866, 0.5);        %e1
				      \coordinate (F) at (-0.866, 0.5);       %e4

				      % Draw all 7 lines of the Fano plane with arrows in middle
				      % Line 1: e3 -> e4 -> e6
				      \draw[thick, midarrow] (B) -- (F);
				      \draw[thick, midarrow] (F) -- (A);
				      % Line 2: e6 -> e1 -> e5
				      \draw[thick, midarrow] (A) -- (E);
				      \draw[thick, midarrow] (E) -- (C);
				      % Line 3: e5 -> e2 -> e3
				      \draw[thick, midarrow] (C) -- (D);
				      \draw[thick, midarrow] (D) -- (B);
				      % Line 4: e6 -> e7 -> e2
				      \draw[thick, midarrow] (A) -- (O);
				      \draw[thick, midarrow] (O) -- (D);
				      % Line 5: e5 -> e7 -> e4
				      \draw[thick, midarrow] (C) -- (O);
				      \draw[thick, midarrow] (O) -- (F);
				      % Line 6: e3 -> e7 -> e1
				      \draw[thick, midarrow] (B) -- (O);
				      \draw[thick, midarrow] (O) -- (E);
				      % Line 7: e1, e2, e7 (as a circle passing through these three points)
				      \draw[thick, decoration={markings, mark=at position 0.166 with {\arrow{<}}, mark=at position 0.5 with {\arrow{<}}, mark=at position 0.833 with {\arrow{<}}}, postaction={decorate}]
				      (0, 0) circle (1);
				      % Place the labels
				      \node at (A) [above]        {\(e_6\)};
				      \node at (B) [below left]   {\(e_3\)};
				      \node at (C) [below right]  {\(e_5\)};
				      \node at (D) [below]        {\(e_7\)};
				      \node at (E) [right]        {\(e_2\)};
				      \node at (F) [left]         {\(e_1\)};
				      \node at (O) [right]        {\(e_4\)};

				      % Add dots for the points
				      \filldraw[black] (A) circle (2pt);
				      \filldraw[black] (B) circle (2pt);
				      \filldraw[black] (C) circle (2pt);
				      \filldraw[black] (D) circle (2pt);
				      \filldraw[black] (E) circle (2pt);
				      \filldraw[black] (F) circle (2pt);
				      \filldraw[black] (O) circle (2pt);
			      \end{tikzpicture}
		      \end{center}
		      If \(e_i \to e_j \to e_k\) is a directed line in the Fano plane, then \(e_i e_j = e_k\), \(e_j e_k = e_i\), and \(e_k e_i = e_j\).
		      The multiplication is anti-commutative, i.e., \(e_i e_j = - e_j e_i\) for all \(i\neq j\).
		      And we have \(e_i^2 = -1\) for all \(i\).
	\end{enumerate}
	\Yang{To be checked...}
\end{remark}

Let \((M,J)\) be an almost complex manifold.
Then the complexified tangent bundle \(TM_\bbC  := TM \otimes_\bbR \bbC\) splits into the direct sum of two complex subbundles
\[ T M_\bbC = T^{1,0}M \oplus T^{0,1}M, \]
where
\[ T^{1,0}M := \ker (\im \id - J), \quad T^{0,1}M := \ker (\im \id + J). \]
We have \(\overline{T^{1,0}M} = T^{0,1}M\) and both \(T^{1,0}M\) and \(T^{0,1}M\) are complex vector bundles of rank \(n\).
This decomposition induces a decomposition of the complexified cotangent bundle
\[ \Omega^1(M) := (TM_\bbC)^* = (T^{1,0}M)^* \oplus (T^{0,1}M)^* =: \Omega^{1,0}(M) \oplus \Omega^{0,1}(M). \]
More generally, for every \(p,q \geq 0\), define
\[ \Omega^{p,q}(M) := \wedge^p (T^{1,0}M)^* \otimes \wedge^q (T^{0,1}M)^* \subset \wedge^{p+q} \Omega^1(M). \]
Then we have the decomposition
\[ \Omega^k(M) \coloneqq \wedge^k \Omega^1(M) = \bigoplus_{p+q=k} \Omega^{p,q}(M). \]
The elements of \(\Omega^{p,q}(M)\) are called \emph{differential forms of type \((p,q)\)} or \emph{\((p,q)\)-forms} for short.
% \Yang{To be continued...}

Recall the \emph{exterior differential operator} \(\upd:\Omega^k(M) \to \Omega^{k+1}(M)\) is locally given by
\[ \upd\left(\sum_I f_I \upd x_I\right) = \sum_I \sum_{j=1}^{2n} \frac{\partial f_I}{\partial x_j} \upd x_j \wedge \upd x_I, \]
where \(I\) runs over all multi-indices with \(|I| = k\) and \(x_1, \ldots, x_{2n}\) are local coordinates on \(M\).

\begin{proposition}\label{prop:decomposition_of_differential_operator_on_almost_complex_manifold}
	There exist differential operators
	\[ \partial: \Omega^{p,q}(M) \to \Omega^{p+1,q}(M), \quad \mu: \Omega^{p,q}(M) \to \Omega^{p+2,q-1}(M) \]
	such that
	\[ \upd = \partial + \overline{\partial} + \mu + \overline{\mu}. \]
\end{proposition}
In a diagram:
\begin{center}
	\begin{tikzpicture}
		\matrix (m) [matrix of math nodes, column sep=2em, row sep=2em] {
			\Omega^{p-1,q+2} & \bullet        & \bullet        & \bullet          & \bullet      \\
			\bullet          & \Omega^{p,q+1} & \bullet        & \bullet          & \bullet      \\
			\bullet          & \Omega^{p,q}   & \Omega^{p+1,q} & \bullet          & \bullet      \\
			\bullet          & \bullet        & \bullet        & \Omega^{p+2,q-1} & \bullet      \\
			                 &                &                & \Omega^{k}       & \Omega^{k+1} \\
		};
		% Draw a horizontal line before the row containing \Omega^k (above row 5)
		\draw (m-4-1.south) -- (m-4-5.south);
		% Draw a dashed line passing through \Omega^{p-1,q+1} and \Omega^k
		\draw[dashed] (m-2-1) -- (m-5-4);
		\draw[dashed] (m-1-1) -- (m-5-5);
		% Draw arrows for the operators
		\draw[->] (m-3-2) -- (m-3-3) node[midway, above] {\(\partial\)};
		\draw[->] (m-3-2) -- (m-2-2) node[midway, right] {\(\overline{\partial}\)};
		\draw[->] (m-3-2) -- (m-4-4) node[midway, right] {\(\mu\)};
		\draw[->] (m-3-2) -- (m-1-1) node[midway, above] {\(\overline{\mu}\)};
		\draw[->] (m-5-4) -- (m-5-5) node[midway, above] {\(\upd\)};
	\end{tikzpicture}
\end{center}
\begin{proof}[Proof of \cref{prop:decomposition_of_differential_operator_on_almost_complex_manifold}]

	\Yang{To be continued...}
\end{proof}

% \begin{proposition}\label{prop:properties_of_differential_operators}
%     The operators \(\partial\) and \(\mu\) satisfy the following properties:
%     \begin{enumerate}
%         \item 
%     \end{enumerate}
%     \Yang{To be continued...}
% \end{proposition}

\begin{definition}\label{def:Nijenhuis_operator}
	The operator \(\mu\) in \cref{prop:decomposition_of_differential_operator_on_almost_complex_manifold} is called the \emph{Nijenhuis operator} of the almost complex structure \(J\).
	If \(\mu = 0\), then \(J\) is called \emph{integrable}.
	In this case, we have \(\upd = \partial + \overline{\partial}\).
	% \Yang{To be continued...}
\end{definition}

\begin{example}\label{eg:Nijenhuis_operator_on_S6_for_the_almost_complex_structure_on_S6_by_octonion}
	Let \(J\) be the almost complex structure on \(S^6\) defined in \cref{eg:almost_complex_structure_on_S6}.

	\Yang{To be continued...}
\end{example}

\begin{proposition}\label{prop:Nijenhuis_operator_on_complex_manifold}
	Let \((M,J)\) be an almost complex manifold.
	If \(J\) is induced by a complex structure on \(M\), then \(J\) is integrable, i.e., the Nijenhuis operator \(\mu = 0\).
\end{proposition}
\begin{proof}
	\Yang{To be continued...}
\end{proof}

The converse of \cref{prop:Nijenhuis_operator_on_complex_manifold} is also true, which is the famous Newlander-Nirenberg theorem.
\Yang{To add reference...}

\begin{theorem}[Newlander-Nirenberg Theorem]\label{thm:newlander_nirenberg_theorem}
	Let \((M,J)\) be an almost complex manifold of real dimension \(2n\).
	If \(\mu = 0\), then \(J\) is induced by a complex structure on \(M\).
\end{theorem}

\begin{proposition}\label{prop:integrable_almost_complex_structure_and_partial_square_zero}
	Let \((M,J)\) be an almost complex manifold.
	Then \(J\) is integrable if and only if \(\partial^2 = 0\).
\end{proposition}



\subsection{Meromorphic functions}

\begin{definition}\label{def:meromorphic_functions}
	Let \(M\) be a complex manifold.
	A \emph{meromorphic function} on \(M\) is a holomorphic map \(f: M \to \bbC\bbP^1\).

	The set of meromorphic functions on \(M\) is denoted by \(\Mer(M)\) or \(\fkk(M)\).
\end{definition}

\begin{proposition}\label{prop:meromorphic_functions_and_holomorphic_functions}
	Let \(M\) be a complex manifold.
	Then there is a natural inclusion \(\Hol(M) \hookrightarrow \Mer(M)\).
	Moreover, we have \(\Mer(M) = \Frac(\Hol(M))\), i.e., every meromorphic function can be expressed as a quotient of two holomorphic functions.
	\Yang{to be checked.}
\end{proposition}

\begin{proposition}\label{prop:meromorphic_functions_is_a_field}
	Let \(M\) be a complex manifold.
	The set of meromorphic functions on \(M\) forms a field under the usual addition and multiplication of functions.
\end{proposition}

\Yang{To be complemented and revised.}



\begin{proposition}\label{prop:schwarz_lemma_in_several_variables}
	Let \(f:U \to \bbC\) be a holomorphic function defined on an open subset \(U \subset \bbC^n\).
	Suppose that \(f\) has order \(k\) at a point \(x \in U\).
	Then there exists a neighborhood \(\overline{B(x,r)} \subset U\) of \(x\) such that
	\[
		|f(z)| \leq C |z - x|^k, \quad \forall z \in \overline{B(x,r)},
	\]
	where \(C = \sup_{z \in \partial \overline{B(x,r)}} \left| f(z) \right|\).
	\Yang{To be revised.}
\end{proposition}

\begin{theorem}[Siegel theorem on function fields]\label{thm:Siegel_theorem_on_function_fields_of_complex_manifolds}
	Let \(X\) be a connected and compact complex manifold of dimension \(n\).
	Then the field of meromorphic functions on \(X\) satisfies
	\[
		\trdeg_{\bbC} \fkk(X) \leq n.
	\]
\end{theorem}
\begin{proof}
	Let \(\{f_1, f_2, \ldots, f_{n+1}\} \subset \fkk(X)\) be meromorphic functions on \(X\).
	We want to find \(P \in \bbC[x_1, x_2, \ldots, x_{n+1}] \setminus \{0\}\) such that
	\[ P(f_1, f_2, \ldots, f_{n+1}) = 0. \]

	\begin{step}
		Let \(z \in X\), there exists \(g_1, g_2, \ldots, g_{n+1}, h \in \Hol(X)\) such that \(f_i = g_i / h\) for each \(1 \leq i \leq n+1\).
	\end{step}

	\Yang{To be revised and complemented.}
\end{proof}

\Yang{To be revised and complemented.}

